\chapter{Testing and CI/CD}

% ============================================================
\section{Overview of Testing Strategy}
% ============================================================
Prism AI Browser is tested at multiple levels corresponding to its layered architecture: the Chromium browser layer (C++), the AI agent backend (Python), the chat/homepage frontend (JavaScript), and the end-to-end integration of all components via CDP. The testing strategy combines manual exploratory testing during development with automated checks embedded in the build pipeline.

% ============================================================
\section{Unit Testing}
% ============================================================

\subsection{Browser Layer (C++)}
Chromium's own GTest-based testing framework is used for unit testing the Prism-specific C++ components:
\begin{itemize}[leftmargin=*]
    \item \textbf{PrismAIButton:} Tests verify that \texttt{OnButtonPressed()} constructs the correct \texttt{OpenURLParams} (URL \texttt{http://127.0.0.1:7788}, disposition \texttt{NEW\_FOREGROUND\_TAB}, transition \texttt{PAGE\_TRANSITION\_TYPED}).
    \item \textbf{CDP Switch Injection:} A unit test confirms that launching the browser without \texttt{--remote-debugging-port} results in the switch being appended with value \texttt{9222}, and that an existing value is not overwritten.
    \item \textbf{Toolbar Integration:} \texttt{ToolbarView} tests assert that the \texttt{prism\_ai\_button\_} child view is non-null after \texttt{Init()} and positioned in the correct order relative to the Home button.
\end{itemize}

\subsection{AI Agent Backend (Python)}
Python's built-in \texttt{unittest} / \texttt{pytest} framework is used to test the AI agent server modules:
\begin{itemize}[leftmargin=*]
    \item \textbf{LLMEngine:} Mock tests verify that prompts are correctly formatted and that the action plan parser produces a well-formed list of \texttt{Action} objects from LLM output.
    \item \textbf{CDPInterface:} Stubbed WebSocket connections test that \texttt{navigate()}, \texttt{evaluate()}, and \texttt{click()} methods produce correctly structured CDP message payloads.
    \item \textbf{BrowserController:} Tests verify that the action dispatcher routes each action type to the correct CDP method and handles errors (timeouts, invalid selectors) gracefully.
\end{itemize}

\subsection{Aura AI Chat Frontend (JavaScript)}
\begin{itemize}[leftmargin=*]
    \item \textbf{TaskExecutor:} Each regex task pattern is tested in isolation against a battery of natural language command variants to confirm correct handler dispatch (e.g., ``open youtube and look for cocomelon'' maps to the YouTube search handler).
    \item \textbf{Exponential Backoff Retry:} Unit tests mock the fetch API to return 503 responses and verify that the retry logic fires exactly 5 times with the expected delays (1 s, 2 s, 4 s, 8 s, 16 s) before surfacing an error.
    \item \textbf{Conversation Storage:} Tests confirm that loading, saving, importing, and exporting conversation sessions from \texttt{localStorage} round-trips correctly.
    \item \textbf{Image Generation:} Mocked OpenRouter API responses are tested to confirm Base64 images are rendered inline and the fallback URL path is triggered when Base64 data is absent.
\end{itemize}

% ============================================================
\section{Integration Testing}
% ============================================================

\subsection{Browser-to-Agent Integration}
\begin{itemize}[leftmargin=*]
    \item The integration test suite starts both the compiled Prism browser and the Python AI agent server, then drives actions through the CDP WebSocket to confirm end-to-end command execution.
    \item A test script submits the command ``Navigate to example.com'' and asserts that the active tab's URL changes to \texttt{https://example.com} within a 5-second timeout.
    \item Tab creation, closure, and switching are similarly exercised end-to-end.
\end{itemize}

\subsection{CDP Availability Test}
A lightweight Python script connects to \texttt{ws://127.0.0.1:9222/json} immediately after browser startup and asserts that a valid JSON list of targets is returned, confirming the CDP always-on feature works correctly.

\subsection{Homepage DevSpace Integration}
The homepage is tested in a locally served environment to confirm:
\begin{itemize}[leftmargin=*]
    \item \texttt{quotes.json} is fetched without CORS errors at \texttt{http://localhost:8000}.
    \item Clock updates in real time and persists the selected format across page reloads.
    \item All Dev Space tool cards open their respective overlays or links correctly.
    \item The Matrix background (\texttt{matrix.js}) renders without JavaScript errors across Chromium and Firefox.
\end{itemize}

% ============================================================
\section{System and Acceptance Testing}
% ============================================================
\begin{itemize}[leftmargin=*]
    \item \textbf{Manual Exploratory Testing:} Each release candidate is tested by team members across Windows and Linux for all major features: AI button, voice commands, image generation, task executor, and homepage.
    \item \textbf{Cross-Platform Build Verify:} The fully patched build is verified to compile without errors on Ubuntu 22.04 (primary CI target) and Windows 11.
    \item \textbf{Privacy Audit:} Network traffic is inspected with \texttt{mitmproxy} during AI feature usage to confirm that no user data or page content is sent to external servers. Only OpenRouter API calls (for image generation) and Gemini/HuggingFace calls (for chat) use the internet, and these are user-initiated with explicit API key configuration.
    \item \textbf{Performance Benchmarking:} Browser startup time is compared against the unpatched Chromium base revision to confirm the CDP switch injection adds $<$50 ms overhead.
\end{itemize}

% ============================================================
\section{CI/CD Pipeline for Prism Browser}
% ============================================================

\subsection{Pipeline Architecture}
The Prism CI/CD pipeline is implemented as a GitHub Actions workflow that validates, builds, and optionally releases a Prism build on every push to the main branch. Figure~\ref{fig:cicd} shows the pipeline stages.

\begin{figure}[H]
\centering
\resizebox{\textwidth}{!}{%
\begin{tikzpicture}[
    stage/.style={draw, rounded corners=5pt, minimum width=3.4cm,
                  minimum height=1.0cm, align=center, fill=blue!10, font=\small},
    dec/.style  ={draw, diamond, aspect=2.5, minimum width=2.8cm,
                  minimum height=1.0cm, align=center, fill=yellow!20, font=\small},
    fail/.style ={draw, rounded corners=5pt, minimum width=3.4cm,
                  minimum height=1.0cm, align=center, fill=red!15, font=\small},
    arr/.style  ={-Stealth, thick},
    lbl/.style  ={font=\scriptsize, fill=white, inner sep=2pt}
]

% -------------------------------------------------------
%  Row 1  (y = 0)  — left → right
% -------------------------------------------------------
\node[stage] (sync)  at (0,    0) {1.\ Chromium Sync\\{\scriptsize(\texttt{gclient sync})}};
\node[stage] (base)  at (5.5,  0) {2.\ Base Revision\\Check};
\node[dec]   (bchk)  at (10.5, 0) {Revision\\match?};
\node[stage] (patch) at (15.5, 0) {3.\ Apply Patches\\{\scriptsize(\texttt{apply\_patches.sh})}};

% -------------------------------------------------------
%  Row 2  (y = -3.8)  — right → left
% -------------------------------------------------------
\node[stage] (build) at (15.5,-3.8) {4.\ GN Gen +\\Ninja Build};
\node[dec]   (bres)  at (10.5,-3.8) {Build\\passes?};
\node[stage] (test)  at (5.5, -3.8) {5.\ Run Unit +\\Integration Tests};
\node[dec]   (tres)  at (1.5, -3.8) {Tests\\pass?};

% -------------------------------------------------------
%  Row 3  (y = -7.8)
% -------------------------------------------------------
\node[stage] (rel)   at (0,   -7.8) {6.\ Package\\Binary};
\node[stage] (pub)   at (5.5, -7.8) {7.\ Publish Release\\Artefact};
\node[fail]  (abrt)  at (10.5,-7.8) {Abort \& Notify\\(email / Slack)};

% -------------------------------------------------------
%  Arrows — Row 1
% -------------------------------------------------------
\draw[arr] (sync)  -- (base);
\draw[arr] (base)  -- (bchk);
\draw[arr] (bchk)  -- node[lbl, above]{Yes} (patch);

% bchk No → re-sync (arc below row 1)
\draw[arr] (bchk.south) -- ++(0,-0.7) -- ++(-12.0,0)
           node[lbl, above, pos=0.5]{No --- re-sync}
           |- (sync.south);

% -------------------------------------------------------
%  Arrows — Row 1 → Row 2
% -------------------------------------------------------
\draw[arr] (patch.south) -- (build.north);

% -------------------------------------------------------
%  Arrows — Row 2
% -------------------------------------------------------
\draw[arr] (build) -- (bres);
\draw[arr] (bres)  -- node[lbl, above]{Yes} (test);
\draw[arr] (test)  -- (tres);

% bres No → abort (straight down, same x)
\draw[arr] (bres.south) -- ++(0,-0.7) -|
           node[lbl, above, pos=0.2]{No} (abrt.north);

% tres No → abort
\draw[arr] (tres.south) -- ++(0,-1.5) -|
           node[lbl, above, pos=0.15]{No} (abrt.south west);

% -------------------------------------------------------
%  Arrows — Row 2 → Row 3
% -------------------------------------------------------
\draw[arr] (tres.west) -- ++(-0.5,0) |-
           node[lbl, right, pos=0.4]{Yes} (rel.north);

% -------------------------------------------------------
%  Arrows — Row 3
% -------------------------------------------------------
\draw[arr] (rel) -- (pub);

\end{tikzpicture}}
\caption{CI/CD Pipeline -- Prism AI Browser}
\label{fig:cicd}
\end{figure}

\subsection{Stage Details}

\begin{longtable}{|p{1.5cm}|p{3.5cm}|p{9cm}|}
\hline
\textbf{Stage} & \textbf{Tool / Script} & \textbf{Description} \\ \hline
\endhead
1 & \texttt{gclient sync} / \texttt{fetch chromium} & Syncs the Chromium source tree to the SHA recorded in \texttt{BASE\_REVISION.txt}. Skippable with \texttt{--skip-sync} flag for speed. \\ \hline
2 & \texttt{apply\_patches.sh} (revision check) & Reads \texttt{BASE\_REVISION.txt} and compares it with \texttt{git rev-parse HEAD} in the Chromium checkout. Fails fast if there is a mismatch. Bypassable with \texttt{--skip-base-check}. \\ \hline
3 & \texttt{apply\_patches.sh} & Applies \texttt{*.patch} files first, then \texttt{*.diff} files, in alphabetical order from the \texttt{patches/} directory. Supports a custom patch directory via \texttt{--patch-dir}. Dry-run mode via \texttt{--check}. \\ \hline
4 & \texttt{build\_prism.sh} $\rightarrow$ \texttt{gn gen} + \texttt{ninja} & Generates the build system with \texttt{gn gen out/Default} and builds the \texttt{chrome} target via \texttt{ninja -C out/Default chrome}. Supports extra GN args via \texttt{--gn-args}. \\ \hline
5 & GTest / pytest & Runs compiled C++ unit tests (\texttt{browser\_tests}, \texttt{unit\_tests}) and Python agent tests. Reports test results to CI. \\ \hline
6 & Custom packaging & Bundles the Prism binary, homepage (\texttt{homepage-devspace/}), and Aura AI chat (\texttt{ai-chat/}) into a distributable archive (\texttt{.tar.gz} on Linux, \texttt{.zip} on Windows). \\ \hline
7 & GitHub Releases API & Uploads the archive as a GitHub Release asset tagged with the Chromium base revision and a Prism version string. \\ \hline
\caption{CI/CD Pipeline Stage Details}
\end{longtable}

\subsection{Key CI/CD Properties}
\begin{itemize}[leftmargin=*]
    \item \textbf{Reproducibility:} Every build is anchored to a specific Chromium revision via \texttt{BASE\_REVISION.txt}, ensuring builds are deterministic and patches apply cleanly.
    \item \textbf{Incremental Builds:} The Ninja build system caches compilation results; only files modified by patches or upstream changes are recompiled, significantly reducing build time on iterative runs.
    \item \textbf{Patch Drift Detection:} If upstream Chromium changes a patched file, the \texttt{git apply} step will fail with a conflict report, immediately alerting contributors to regenerate affected patches.
    \item \textbf{Separation of Concerns:} Feature patches are kept in separate, numbered files (\texttt{0001-*}, \texttt{0002-*}, \ldots), enabling individual features to be reviewed, reverted, or cherry-picked independently.
    \item \textbf{Minimal Repository Size:} The public Prism repository stores only patch files, scripts, and documentation -- no Chromium source code -- keeping repository size under 1 MB.
\end{itemize}
